\section{Fundamentação e estado da arte}
\label{sec:fundamentacao}

\subsection{O CyBOK e o ciclo de vida de desenvolvimento seguro}

O CyBOK organiza o conhecimento de cibersegurança em áreas de conhecimento agrupadas em categorias \parencite{cybok2021}. O problema deste trabalho se situa na categoria de Segurança de Software e Plataforma, na área de Ciclo de Vida de Desenvolvimento Seguro (Secure Software Lifecycle) \parencite{williams2021}, com apoio da área de Segurança de Software.

Essa área parte da premissa de que segurança não é uma funcionalidade acrescentada ao final, e sim uma propriedade construída e verificada continuamente ao longo do desenvolvimento. Suas características centrais são a modelagem de ameaças no início, a incorporação de práticas seguras no projeto e na codificação, o teste de segurança como atividade recorrente e a melhoria iterativa a cada ciclo. Ela se distingue da área vizinha de Segurança de Software, que foca em categorias concretas de falhas de código, por tratar do processo que decide quais controles aplicar e quando. Distingue-se das dimensões de Segurança de Sistemas e de Infraestrutura por não ser sobre um mecanismo isolado (criptografia, hardware, rede), mas sobre como conduzir o desenvolvimento. E separa-se da dimensão de Ataques e Defesas por atuar na fase de construção, e não na operação e na resposta a incidentes.

A área descreve processos prescritivos de ciclo de vida seguro, entre eles o Security Development Lifecycle (SDL) da Microsoft \parencite{howard2006}, os pontos de contato (Touchpoints) de \textcite{mcgraw2006} e as práticas fundamentais da SAFECode \parencite{safecode2018}, além das adaptações desses processos a contextos como o desenvolvimento ágil e o DevOps. Práticas comuns a esses processos, como modelagem de ameaças, casos de abuso, requisitos de segurança, métricas, padrões de criptografia, gestão de componentes de terceiros, análise estática e teste de penetração, aparecem ao longo deste trabalho e são apontadas onde entram. Três outras áreas do CyBOK dão apoio pontual: Segurança de Software, para as fraquezas do binário; Criptografia, para a licença assinada; e Autenticação, Autorização e Responsabilização, para o licenciamento.

\subsection{O atacante Man-At-The-End}

No modelo MATE, o atacante tem acesso físico e administrativo ao dispositivo em que o software executa \parencite{collberg2009}. Ele pode ler o binário, executá-lo sob um depurador, alterar a memória do processo, interceptar chamadas e reescrever o código. Contra esse atacante, nenhuma proteção que resida inteiramente no cliente é absoluta: com esforço suficiente, quem controla a máquina adultera o binário ou reconstrói o que ele faz. O que as defesas fazem é encarecer o ataque. Por isso a pergunta útil deixa de ser \enquote{o software é seguro?} e passa a ser \enquote{quanto custa atacá-lo, e esse custo supera o valor do que se protege?}.

\subsection{Proteção de software contra o MATE}

A proteção de software contra atacantes MATE conta com corpos de conhecimento consolidados. O projeto europeu ASPIRE organiza as defesas em cinco linhas complementares, ocultação de dados, ocultação de algoritmo, anti-adulteração, atestação remota e renovabilidade, e avalia a proteção pelo custo do ataque, não por um veredito binário de seguro ou inseguro \parencite{aspire}. O padrão OWASP MASVS, embora voltado a aplicativos móveis, traz uma categoria dedicada à resiliência contra engenharia reversa e adulteração, a MASVS-RESILIENCE, usada aqui por analogia: RESILIENCE-1 para integridade da plataforma, RESILIENCE-2 para anti-adulteração, RESILIENCE-3 para dificultar a análise estática e RESILIENCE-4 para dificultar a análise dinâmica \parencite{masvs}. As fraquezas são nomeadas pelo catálogo CWE \parencite{cwe}.

Uma lição central desse estado da arte orienta as expectativas: em um cenário MATE, ocultar um segredo que precisa executar no cliente é, no fim, segurança por obscuridade (CWE-656). A ofuscação e a compilação nativa encarecem a leitura, mas não a impedem. A mitigação que de fato derrota o MATE é não enviar o segredo ao cliente.

\subsection{Métodos de engenharia de segurança}

A condução do processo se apoia em métodos estabelecidos. A modelagem de ameaças usa o STRIDE \parencite{shostack2014}, que classifica ameaças em falsificação de identidade (Spoofing), adulteração (Tampering), repúdio (Repudiation), vazamento de informação (Information disclosure), negação de serviço (Denial of service) e elevação de privilégio (Elevation of privilege). Os casos de abuso descrevem o comportamento do sistema sob um ator malicioso, o oposto das histórias de usuário \parencite{sindre2005,mcgraw2006}. As árvores de ataque organizam os caminhos que um atacante pode percorrer até um objetivo, com nós OR, em que basta um caminho, e AND, em que todos são necessários \parencite{schneier1999}. O Test-Driven Security leva para a segurança a lógica do desenvolvimento guiado por testes: o controle nasce junto com um teste que prova que ele funciona e que, dali em diante, impede a volta da falha \parencite{vehent2018}. O purple teaming junta o time que ataca e o time que defende em torno do mesmo resultado, a cada ciclo.

\subsection{Princípios de projeto seguro}

O CyBOK retoma os princípios de projeto de \textcite{saltzer1975}. Três deles explicam boa parte das decisões e do resultado deste trabalho. O padrão seguro em caso de falha (fail-safe defaults) orientou as respostas de integridade, que fecham o treino quando falta qualquer evidência de que o pacote está íntegro. A separação de privilégio orientou as duas guardas independentes e a chave privada do fornecedor fora do cliente. E o projeto aberto (open design), segundo o qual a segurança não deve depender do segredo do mecanismo, antecipa a diferença entre as duas histórias: a licença não depende de esconder nada no cliente, e fechou; o algoritmo depende, e só encareceu.
